\section*{\SGAInlineRefTarget{sga3:x:section:4}4. Case of an arbitrary base. Quasi-isotriviality theorem}
\phantomsection
\addcontentsline{toc}{section}{4. Case of an arbitrary base. Quasi-isotriviality theorem}

Let \(A\) be a local ring. Recall, following Nagata, that \(A\) is said to be
\emph{henselian} if every algebra \(B\) finite over \(A\) is a product of local
algebras \(B_i\) finite over \(A\).

\par\medskip\noindent
\textbf{\SGAInlineRefTarget{sga3:x:paragraph:4-0}Recall 4.0.}\footnote{\textit{Editors' note.} This recall has been
spelled out, and \SGARefLink{sga3:x:remark:4-0-1}{Remark 4.0.1} has been added.}\enspace
\emph{Let \(A\) be a henselian local ring, \(k\) its residue field,
\(S=\operatorname{Spec}(A)\), \(S_0=\operatorname{Spec}(k)\), and \(\xi\) a
geometric point of \(S_0\). Then the functor}
\[
  X\longmapsto X_0=X\times_S S_0
\]
\emph{is an equivalence between the category of étale coverings of \(S\) and
the analogous category over \(S_0\) (cf.\ EGA IV\(_4\), \S~18.5).
Consequently (cf.\ SGA 1, V),}
\[
  \pi_1(S_0,\xi)=\pi_1(S,\xi).
\]

Suppose in addition that \(A\) is noetherian, and write \(A'\) for its
completion and \(S'=\operatorname{Spec}(A')\). Then \(A'\) is a complete
noetherian local ring, hence henselian (\emph{loc.\ cit.}, 18.5.14), and the
functor
\[
  X\longmapsto X'=X\times_S S'
\]
from the category of étale coverings of \(S\) to the analogous category over
\(S'\) fits into a commutative diagram
\[
  {\small
\begin{tikzcd}[column sep=2.5em, row sep=2.8em]
  \mathrm{\acute{E}tCov}(S)
    \arrow[rr]
    \arrow[dr,"\sim"']
  &&
  \mathrm{\acute{E}tCov}(S')
    \arrow[dl,"\sim"]
  \\
  &
  \mathrm{\acute{E}tCov}(S_0)
  &
\end{tikzcd}
}

\]
and is therefore also an equivalence of categories, whence
\(\pi_1(S_0,\xi)=\pi_1(S',\xi)\).

\par\medskip\noindent
\textbf{\SGAInlineRefTarget{sga3:x:remark:4-0-1}Remark 4.0.1.}\enspace
Since \(S\) is connected (\(A\) being local), it follows from \SGARefLink{sga3:x:corollary:1:2}{1.2} that the
category of isotrivial groups of multiplicative type over \(S\) is equivalent
to the analogous category over \(S_0\) (and also over \(S'\) if \(A\) is
noetherian).

\par\medskip\noindent
\textbf{\SGAInlineRefTarget{sga3:x:lemma:4-1}Lemma 4.1.}\enspace
\begingroup\itshape
Let \(A\) be a henselian local ring with residue field \(k\),
\(S=\operatorname{Spec}(A)\), and \(S_0=\operatorname{Spec}(k)\).

\smallskip
\noindent
\(\mathrm{(i)}\) The functor
\[
  H\longmapsto H_0=H\times_S S_0
\]
is an equivalence between the category of finite groups of multiplicative type
over \(S\) and the analogous category over \(S_0\).

\smallskip
\noindent
\(\mathrm{(ii)}\) If, in addition, \(A\) is noetherian, writing \(A'\) for its
completion and \(S'=\operatorname{Spec}(A')\), the functor
\[
  H\longmapsto H'=H\times_S S'
\]
is an equivalence between the category of finite groups of multiplicative type
over \(S\) and the analogous category over \(S'\).
\par\endgroup

As in \SGARefLink{sga3:x:paragraph:4-0}{4.0}, the second assertion follows from the first; let us prove the first.
We already know that the functor under consideration is essentially
surjective: every group of multiplicative type \(H_0\) over
\(S_0=\operatorname{Spec}(k)\), being finite and hence of finite type over
\(k\), is isotrivial by \SGARefLink{sga3:x:proposition:1:4}{1.4}, and therefore comes from an isotrivial group of
multiplicative type over \(S\), by \SGARefLink{sga3:x:remark:4-0-1}{4.0.1}.

It remains to prove full faithfulness, that is, that for two finite groups
\(G,H\) of multiplicative type over \(S\), the following map is bijective:
\[
  \operatorname{Hom}_{S\text{-gr.}}(G,H)
  \longrightarrow
  \operatorname{Hom}_{S_0\text{-gr.}}(G_0,H_0).
\]
In other words,\footnote{\textit{Editors' note.} What follows has been added.}
if \(F\) denotes the \(S\)-functor
\(\operatorname{Hom}_{S\text{-gr.}}(G,H)\), the natural map
\[
  \operatorname{Hom}_S(S,F)
  \longrightarrow
  \operatorname{Hom}_{S_0}(S_0,F_0)
\]
induced by the base change \(S_0\to S\) is bijective. In view of Recall \SGARefLink{sga3:x:paragraph:4-0}{4.0},
it is enough to prove the following lemma.

\par\medskip\noindent
\textbf{\SGAInlineRefTarget{sga3:x:lemma:4-2}Lemma 4.2.}\enspace
\emph{Let \(G,H\) be two finite groups of multiplicative type over a
prescheme \(S\). Then
\(F=\operatorname{Hom}_{S\text{-gr.}}(G,H)\) is representable by a finite
étale scheme over \(S\).}

\footnote{\textit{Editors' note.} The original has been expanded in what
follows.}
Let \(f:S'\to S\) be a faithfully flat quasi-compact morphism such that
\(G'\) and \(H'\) are diagonalizable. It is enough to show that \(F_{S'}\) is
represented by a prescheme \(X'\) étale and finite (hence affine) over \(S'\).
Indeed, the descent datum on \(X'\) relative to \(f\) (cf.\ VIII 1.7.2) is then
effective (by SGA 1, VIII 2.1), so there is an \(S\)-prescheme \(X\) such that
\(X\times_S S'=X'\), representing \(F\), and \(X\) is étale and finite over
\(S\) (cf.\ SGA 1, V 5.7 and EGA IV\(_4\), 17.7.3 (ii)).

We may therefore suppose that \(G=D_S(M)\) and \(H=D_S(N)\), where \(M\) and
\(N\) are finite commutative groups (cf.\ \SGARefLink{sga3:viii:proposition:2:1}{VIII 2.1} c)). Then
\(K=\operatorname{Hom}_{\mathrm{gr.}}(N,M)\) is a finite abelian group and,
by \SGARefLink{sga3:viii:corollary:1:5}{VIII 1.5}, there is an isomorphism
\[
  \operatorname{Hom}_{S\text{-gr.}}(G,H)\simeq K_S,
\]
which completes the proof of \SGARefLink{sga3:x:lemma:4-2}{4.2} and hence of \SGARefLink{sga3:x:lemma:4-1}{4.1}.

\par\medskip\noindent
\textbf{\SGAInlineRefTarget{sga3:x:proposition:4-3-0}Proposition 4.3.0.}\footnote{\textit{Editors' note.} This recall,
which is used in the proof of 4.3, has been added (the reference to
EGA IV\(_4\), 18.5.11 was not entirely satisfactory).}\enspace
\begingroup\itshape
Let \(S\) be a henselian local scheme, \(s\) its closed point,
\(g:X\to S\) a morphism locally of finite type, and \(x\) an isolated point
of the fiber \(X_s\).

\smallskip
\noindent
\(\mathrm{(i)}\) Then \(\mathcal O_{X,x}\) is finite over
\(\mathcal O_{S,s}\). (In particular, if the residue-field extension
\(\kappa(x)/\kappa(s)\) is trivial, then
\(\mathcal O_{S,s}\to\mathcal O_{X,x}\) is surjective, by Nakayama's lemma.)

\smallskip
\noindent
\(\mathrm{(ii)}\) If \(g\) is also separated, then
\(X'=\operatorname{Spec}(\mathcal O_{X,x})\) is an open and closed subscheme
of \(X\); that is, there is a decomposition into a disjoint sum
\(X=X'\amalg X''\).
\par\endgroup

\par\smallskip\noindent
\emph{Proof.}
By the local form of Zariski's Main Theorem found in [Pes66], or in
[Ray70, Ch.~IV, Th.~1], \(x\) has an affine open neighborhood
\(U=\operatorname{Spec}(B)\), of finite type and quasi-finite over
\(A=\mathcal O_{S,s}\), and there is an open immersion
\(U\hookrightarrow Y=\operatorname{Spec}(C)\), where \(C\) is a finite
\(A\)-algebra. (N.B. To obtain this conclusion one may also use Chevalley's
semicontinuity theorem (EGA IV\(_3\), 13.1.4), followed by the form of
Zariski's Main Theorem given in \emph{loc.\ cit.}, 8.12.8.)

Since \(A\) is henselian, \(Y\) is the disjoint sum of local schemes
\(Y_1,\ldots,Y_n\), each finite over \(S\), and the points of \(Y\) above
\(s\) are their closed points \(y_1,\ldots,y_n\). Thus \(x=y_i\) for some
\(i\), and
\(\mathcal O_{X,x}=\mathcal O_{U,x}=C_i\) is finite over \(A\). Moreover,
\(X'=\operatorname{Spec}(C_i)\) is an open subscheme of \(U\), hence of
\(X\).

Suppose in addition that \(g\) is separated. Since \(X'\to S\) is finite
(\(C_i\) being finite over \(A\)), so is the immersion \(X'\to X\)
(cf.\ EGA II, 6.1.5 (v)); therefore \(X'\) is also closed in \(X\).
\hfill\(\square\)

\par\medskip\noindent
\textbf{\SGAInlineRefTarget{sga3:x:lemma:4-3}Lemma 4.3.}\enspace
\begingroup\itshape
Let \(A\) be a noetherian henselian local ring, \(A'\) its completion,
\(S=\operatorname{Spec}(A)\), \(S'=\operatorname{Spec}(A')\), and \(s\) the
closed point of \(S\). Let \(G,H\) be two \(S\)-preschemes in groups, with
\(G\) of multiplicative type and of finite type over \(S\), and with \(H\)
locally of finite type and separated over \(S\), \(H_s\) of multiplicative
type, and \(H\) flat over \(S\) at the points of \(H_s\).

Let \(G',H'\) be obtained from \(G,H\) by the base change \(S'\to S\). Then
the following natural map is bijective:
\par\endgroup
\[
  \operatorname{Hom}_{S\text{-gr.}}(G,H)
  \xrightarrow{\ \sim\ }
  \operatorname{Hom}_{S'\text{-gr.}}(G',H').
\]

Since \(S'\) is faithfully flat and quasi-compact over \(S\), SGA 1,
VIII 5.2 shows that this map identifies the first member with the subset of
the second formed by the \(u':G'\to H'\) for which the two inverse images
\(u''_1,u''_2:G''\to H''\) of \(u'\) on
\(S''=S'\times_S S'\), under the projections
\(\operatorname{pr}_1,\operatorname{pr}_2:S''\to S'\), are equal.

It therefore remains to prove that \(u''_1=u''_2\) for every homomorphism of
\(S'\)-groups \(u':G'\to H'\). By the density theorem \SGARefLink{sga3:ix:corollary:4:8}{IX 4.8}, it is enough
to prove that \(u''_1\) and \(u''_2\) agree on \({}_nG''\) for every integer
\(n>0\). (N.B. Here the density theorem is needed essentially in a case
where the base prescheme, namely \(S''\), is not locally noetherian.)

Replacing \(G\) by \({}_nG\), this reduces us to the case in which
\(n\cdot\operatorname{id}_G=0\) for some integer \(n>0\), and hence \(G\)
is finite over \(S\). Likewise let \({}_nH\) be the kernel of the \(n\)-th
power morphism \(\varphi_n\) on \(H\). (N.B. We have not assumed \(H\)
commutative, so \(\varphi_n\), respectively \({}_nH\), need not be a group
homomorphism, respectively a subgroup of \(H\).) Since \({}_nH\) is defined
as the fiber product of \(\varphi_n:H\to H\) and the unit section
\(\varepsilon:S\to H\), its formation commutes with every base change
\(T\to S\):
\[
  ({}_nH)_T={}_n(H_T).
\]

\footnote{\textit{Editors' note.} The original has been expanded in what
follows, on the basis of indications by M.\ Raynaud.}
We use a right subscript \(m\) for reduction modulo
\(\mathfrak m^{m+1}\), where \(\mathfrak m\) is the maximal ideal of \(A\).
Then \(H_m\) is flat over \(S_m\) (because \(H\) is flat over \(S\) at the
points of \(H_0\)); hence, by \SGARefLink{sga3:x:corollary:2:4}{2.4}, \(H_m\) is of multiplicative type and of
finite type, since \(H_0\) is. Thus every
\(({}_nH)_m={}_n(H_m)\) is a subgroup of \(H_m\), of multiplicative type,
finite and flat over \(S_m\).

In particular, \(({}_nH)_0={}_n(H_0)\) is finite over \(S_0\). Since \(S\)
is henselian local and \({}_nH\), like \(H\), is locally of finite type and
separated over \(S\), \SGARefLink{sga3:x:proposition:4-3-0}{Proposition 4.3.0} shows that the local rings of
\({}_nH\) at the points of \(({}_nH)_0\) are finite over \(S\), and that
there is a decomposition as a disjoint sum of open subsets
\[
  \tag{*}\label{sga3:x:lemma:4-3:equation:star}
  {}_nH={}_nH^+\amalg{}_nH^-,
\]
where \(Z={}_nH^+\) is finite over \(S\), while \({}_nH^-\) lies above
\(S-\{s\}\).

Observe that for every finite \(S\)-scheme \(Y\) (such as \(G\)), every
\(S\)-morphism \(Y\to{}_nH\) factors through \(Z\), and that the formation
of the decomposition \((*)\) commutes with the base change \(S'\to S\),
where \(S'=\operatorname{Spec}(A')\) and \(A'\) is the completion of \(A\).

Then \(Z'=Z\times_S S'\) is finite over \(S'\), as is
\(P'=Z'\times_{S'}Z'\). Let \(\nu\) denote the restriction to \(P'\) of
the multiplication on \(H'\), and let \(\sigma\) be the automorphism of
\(P'\) interchanging the two factors. Since \(P'\) is finite over \(S'\),
while \(H'\) is separated and locally of finite type over \(S'\), EGA II
5.4.3 and IV\(_1\) 1.1.3 show that \(Y=\nu(P')\) is a closed subscheme of
\(H'\), universally closed and quasi-compact, hence of finite type and
therefore proper over \(S'\). Moreover, \(Y\to S'\) has finite fibers
(because \(P'\to S'\) does). Since \(S'\) is noetherian, \(Y\to S'\) is
therefore finite (cf.\ EGA III, 4.4.2). Since \(Z'\subseteq Y\) and
\(Z'_m=Y_m\) for every \(m\), Lemma IX 5.0 gives \(Z'=Y\); thus \(Z'\) is
a subgroup of \(H'\). Similarly, the kernel
\(K=\operatorname{Ker}(\nu,\nu\circ\sigma)\) is a closed subscheme of
\(P'\) such that \(K_m=P'_m\) for every \(m\) (because
\(Z'_m={}_nH_m\) is commutative). Hence \(K=P'\), that is, \(Z'\) is a
commutative subgroup of \(H'\).

On the other hand, every reduction \(Z'_m={}_nH_m\) is flat over \(S_m\).
The ``local criterion for flatness'' (cf.\ EGA \(0_{\mathrm{III}}\),
10.2.2, or [BAC], III \S~5, Example 1 and Th.~1) therefore shows that
\(Z'\) is flat over \(S'\). Moreover,
\(Z'_0={}_nH_0\) is a finite group of multiplicative type over \(S_0\),
and hence isotrivial by \SGARefLink{sga3:x:proposition:1:4}{1.4}. It follows from \SGARefLink{sga3:x:corollary:3:8}{3.8} b) that \(Z'\) is of
multiplicative type (and isotrivial) over \(S'\). Since \(S'\to S\) is
faithfully flat and quasi-compact, the multiplication
\(Z\times_S Z\to H\) factors through \(Z\) and makes \(Z\) a subgroup of
\(H\), finite over \(S\) and of multiplicative type (because \(Z'\) is).

Finally, by the observations made above concerning \((*)\), \(Z'\) is the
``finite part'' of \({}_nH'\), and therefore \(u':G'\to{}_nH'\) takes its
values in \(Z'\). Since \(Z\), as well as \(G={}_nG\), is finite over \(S\)
and of multiplicative type, \SGARefLink{sga3:x:lemma:4-1}{Lemma 4.1} shows that \(u'\) comes from a
\(u:G\to Z\), and hence \(u''_1=u''_2\). This completes the proof of \SGARefLink{sga3:x:lemma:4-3}{4.3}.

\par\medskip\noindent
\textbf{\SGAInlineRefTarget{sga3:x:lemma:4-4-0}Lemma 4.4.0.}\footnote{\textit{Editors' note.} This lemma, used
several times below, has been added.}\enspace
\begingroup\itshape
Let \(A\) be a ring, \(S=\operatorname{Spec}(A)\), and let \(H\) be
an \(S\)-group scheme of multiplicative type and of finite type,
quasi-isotrivial (respectively isotrivial). Let \((A_i)\) be a filtered
family of subrings of \(A\) whose inductive limit is \(A\), and put
\(S_i=\operatorname{Spec}(A_i)\).

Then there are an index \(i\) and an \(S_i\)-group scheme \(H_i\), of
multiplicative type and of finite type, quasi-isotrivial (respectively
isotrivial), such that \(H=H_i\times_{S_i}S\).
\par\endgroup

\par\medskip\noindent
\textbf{\SGAInlineRefTarget{sga3:x:theorem:4-4}Theorem 4.4.}\enspace
\begingroup\itshape
Let \(S\) be a prescheme, let \(H\) be an affine \(S\)-prescheme in
groups of finite presentation over \(S\), and let \(s\in S\). Suppose:

\smallskip
\noindent
\(\alpha)\) \(H\) is flat over \(S\) at the points of \(H_s\);

\smallskip
\noindent
\(\beta)\) \(H_s\) is of multiplicative type.

\smallskip
Then there exist an étale morphism \(S'\to S\), a point \(s'\in S'\) above
\(s\) such that \(\kappa(s')/\kappa(s)\) is trivial, and an open and closed
subgroup \(G'\) of \(H'=H\times_S S'\), of multiplicative type and finite
type, and isotrivial, such that \(G'_{s'}=H'_{s'}\).
\par\endgroup

\footnote{\textit{Editors' note.} The reductions that follow have been
spelled out (the original said: ``One reduces at once to the case where
\(S\) is the spectrum of a noetherian local ring \(A\), and where \(s\) is
the closed point of \(S\).'')}
\noindent\textbf{a)}
Write provisionally \(T=\operatorname{Spec}(\mathcal O_{S,s})\), and let us
show, by ``descent of the conclusions'', that one may reduce to the case
\(S=T\). Suppose indeed that we have found an étale morphism \(g:T'\to T\),
a point \(s'\in T'\), and a subgroup \(G'\) of \(H'=H_{T'}\) satisfying the
conditions of the statement. Replacing \(T'\) by an affine open
neighborhood of \(s'\), we may assume that \(T'\) is of finite presentation
over \(T\).

Since \(G'\) is isotrivial, there is a finite étale surjective morphism
\(f:\widetilde T\to T'\) such that
\(\widetilde G=G'\times_{T'}\widetilde T\) is isomorphic to
\(D_{\widetilde T}(M)\), where \(M\) is a finitely generated abelian group.
Since \(T'\), \(H\), \(\widetilde T\), and \(G'\) are of finite presentation
over \(T=\operatorname{Spec}(\mathcal O_{S,s})\), and
\(\mathcal O_{S,s}\) is the inductive limit of the subalgebras
\(A_i=\mathcal O_S(S_i)\), where \(S_i\) runs through the affine open
neighborhoods of \(s\) in \(S\), EGA IV\(_3\), 8.8.2 (and Exposé
VI\(_{\mathrm B}\), 10.2 and 10.3) gives an index \(i\), \(S_i\)-preschemes
\(S'_i,\widetilde T_i\) of finite presentation (respectively an
\(S_i\)-prescheme in groups \(G'_i\) of finite presentation), and morphisms
\[
  g_i:S'_i\to S_i,\qquad f_i:\widetilde T_i\to S'_i,
  \qquad
  u_i:G'_i\to H\times_S S'_i,
\]
from which \(g,f\), respectively \(u:G'\to H'\), are obtained by base
change. Taking \(i\) large enough, \(g_i\) is étale, \(f_i\) is finite
étale surjective, and \(u_i\) is an open and closed immersion
(cf.\ EGA IV, 8.10.5 and 17.7.8).

Now \(\widetilde G\) comes from both
\(\widetilde G_i=G'_i\times_{S'_i}\widetilde T_i\) and
\(D_{\widetilde T_i}(M)\). By EGA IV\(_3\), 8.8.2 (i) (and
VI\(_{\mathrm B}\), 10.2), there is an index \(j\geq i\) such that
\(\widetilde G_j\simeq D_{\widetilde T_j}(M)\), so \(G'_j\) is an
isotrivial group of multiplicative type. Let \(s'_j\) be the image of
\(s'\) in \(S'_j\). Then the étale morphism
\(S'_j\to S_j\hookrightarrow S\), the point \(s'_j\), and the open and
closed subgroup \(G'_j\) of \(H\times_S S'_j\) satisfy the conditions of
the statement. Thus we may assume that \(S=\operatorname{Spec}(A)\) is
local, with closed point \(s\).

\par\smallskip\noindent
\textbf{b)}
The ring \(A\) is the inductive limit of local rings \(A_i\) which are
localizations of finitely generated \(\mathbb Z\)-algebras; put
\(S_i=\operatorname{Spec}(A_i)\). Let us show that hypotheses \(\alpha\)
and \(\beta\) ``descend'' to some \(A_i\). Since \(H\to S\) is of finite
presentation, the set of \(x\in H\) at which \(H\) is flat over \(S\) is
an open set \(W\). It contains \(H_s\) by hypothesis, and hence contains a
quasi-compact open set \(V\) containing \(H_s\) (since the affine scheme
\(H_s\) can be covered by finitely many affine open subsets contained in
\(W\)). The open immersion \(\tau:V\hookrightarrow H\) is of finite
presentation, so \(V\to S\) is also of finite presentation.

By EGA IV\(_3\), 8.8.2 (and Exposé VI\(_{\mathrm B}\), 10.2 and 10.3),
there are an index \(i\), an \(S_i\)-prescheme \(V_i\) of finite
presentation (respectively an \(S_i\)-prescheme in groups \(H_i\) of
finite presentation), and an \(S_i\)-morphism \(\tau_i:V_i\to H_i\), from
which \(V,H,\tau\) are obtained by base change \(S\to S_i\). Taking \(i\)
large enough, \(\tau_i\) is an open immersion and \(V_i\) is flat over
\(S_i\), by EGA IV\(_3\), 8.10.5 and 11.2.6. Let \(s_i\) be the image of
\(s\) in \(S_i\). Since the open immersion
\((V_i)_{s_i}\hookrightarrow(H_i)_{s_i}\), after base change
\(\kappa(s_i)\to\kappa(s)\), gives the equality \(V_s=H_s\), one already
has \((V_i)_{s_i}=(H_i)_{s_i}\). Thus \(H_i\) is flat over \(S_i\) at the
points of \((H_i)_{s_i}\). Finally, \((H_i)_{s_i}\) is of multiplicative
type, because
\[
  H_s=(H_i)_{s_i}\otimes_{\kappa(s_i)}\kappa(s)
\]
is. Hence \((S_i,H_i,s_i)\) satisfies the hypotheses of \SGARefLink{sga3:x:theorem:4-4}{4.4}, and if the
desired assertion holds for this triple, it also holds for \((S,H,s)\) by
base change. We are reduced to the case in which \(A\) is local and
noetherian. We now distinguish two cases.

\par\smallskip\noindent
\(\mathbf{1^\circ}\) \emph{\(A\) is local noetherian and henselian.}
Let \(\widehat A\) be its completion,
\(\widehat S=\operatorname{Spec}(\widehat A)\), and
\(\widehat H=H\times_S\widehat S\). Applying \SGARefLink{sga3:x:theorem:3:7}{Theorem 3.7}, we find an
isotrivial \(\widehat S\)-group \(\widehat G\) of multiplicative type and
of finite type, together with a homomorphism
\[
  \widehat u:\widehat G\longrightarrow\widehat H
\]
which is both an open and a closed immersion, and which induces an
isomorphism
\[
  \widehat u_0:\widehat G_0\xrightarrow{\ \sim\ }\widehat H_0.
\]
By \SGARefLink{sga3:x:remark:4-0-1}{Remark 4.0.1}, base change along \(\widehat S\to S\) induces an
equivalence between the categories of isotrivial groups of multiplicative
type over \(S\) and over \(\widehat S\). In particular, \(\widehat G\)
comes from an isotrivial \(S\)-group \(G\) of multiplicative type and of
finite type. By \SGARefLink{sga3:x:lemma:4-3}{4.3}, \(\widehat u\) comes from a homomorphism
\[
  u:G\longrightarrow H.
\]
Moreover, \(u\) is an open and closed immersion and induces an isomorphism
\(u_0:G_0\to H_0\), since this is true after the faithfully flat
quasi-compact base change \(\widehat S\to S\). This proves \SGARefLink{sga3:x:theorem:4-4}{4.4} in this case
(taking \(S'=S\) and \(s'=s\) in its conclusion).

\par\smallskip\noindent
\(\mathbf{2^\circ}\) \emph{\(A\) is local noetherian.}
The reduction to case \(1^\circ\) is immediate, by applying that case to
the henselization \(A^h\) of \(A\). More precisely, one sees easily (using
SGA 1, I \S~5%
\begingroup
\renewcommand{\thefootnote}{\(\ast\)}
\footnote{Or EGA IV\(_4\), \S~18.6.}
\addtocounter{footnote}{-1}
\endgroup)
that the local rings \(\mathcal O_{S',s'}\) of étale \(S\)-preschemes
\(S'\), equipped with a point \(s'\) above \(s\) for which
\(\kappa(s')/\kappa(s)\) is trivial, form a filtered inductive system whose
inductive limit is a noetherian henselian local ring \(A^h\) (called the
henselization of \(A\)); for details of this construction, due to Nagata in
the normal case, see SGA 4, VIII \S~4. The standard arguments of
EGA IV\(_3\), \S~8 then allow one, as in part a), to deduce from the known
result over the inductive limit \(A^h\) of the \(\mathcal O_{S',s'}\) an
analogous result over one of the pairs \((S',s')\). This proves \SGARefLink{sga3:x:theorem:4-4}{4.4}.

\par\medskip\noindent
\textbf{\SGAInlineRefTarget{sga3:x:corollary:4-5}Corollary 4.5.}\enspace
\emph{Let \(S\) be a prescheme and \(H\) an \(S\)-prescheme in groups of
multiplicative type and of finite type. Then \(H\) is quasi-isotrivial;
that is, \(H\) is split by an étale surjective morphism \(S'\to S\).}

Indeed, let \(s\in S\). By \SGARefLink{sga3:x:theorem:4-4}{4.4}, there are an étale morphism \(S'\to S\), a
point \(s'\in S'\) above \(s\) with \(\kappa(s)=\kappa(s')\), and an
isotrivial subgroup \(G'\) of \(H'\), of multiplicative type and of finite
type, such that \(G'_{s'}=H'_{s'}\). Since \(G'\) and \(H'\) are of
multiplicative type and finite type, \SGARefLink{sga3:ix:corollary:2:9}{IX 2.9} gives an open neighborhood
\(U'\) of \(s'\) such that \(G'|_{U'}=H'|_{U'}\).

\footnote{\textit{Editors' note.} The original has been expanded in what
precedes and what follows.}
Suppose in addition that \(S\) is henselian local, with closed point \(s\).
By EGA IV\(_4\), 18.5.11 b), there is a section \(\sigma\) of \(S'\to S\)
such that \(\sigma(s)=s'\). (N.B. One can see directly that
\(B=\mathcal O_{S',s'}\) equals \(A=\mathcal O_{S,s}\) as follows. By
4.3.0 (i), \(B\simeq A/I\). Since \(B\) is a flat \(A\)-algebra of finite
presentation, \(I\) is a finitely generated ideal (cf.\ EGA IV\(_1\),
1.4.7), and \(I=I\mathfrak m\), where \(\mathfrak m\) is the maximal ideal
of \(A\); hence \(I=0\).) Thus \(H\) is already isotrivial. We obtain:

\par\medskip\noindent
\textbf{\SGAInlineRefTarget{sga3:x:corollary:4-6}Corollary 4.6.}\enspace
\begingroup\itshape
Let \(A\) be a henselian local ring, \(k\) its residue field, and
\(\pi\) the topological Galois group of an algebraic closure of \(k\).

\smallskip
\noindent
\(\mathrm{(i)}\) Every group of multiplicative type and of finite type over
\(S=\operatorname{Spec}(A)\) is isotrivial.

\smallskip
\noindent
\(\mathrm{(ii)}\) The category of these groups over \(S\) is equivalent,
through the functor \(H\mapsto H\otimes_A k\), to the analogous category
over \(S_0=\operatorname{Spec}(k)\). It is therefore anti-equivalent, by
\SGARefLink{sga3:x:proposition:1:4}{1.4}, to the category of Galois modules under \(\pi\) which are finitely
generated as \(\mathbb Z\)-modules.
\par\endgroup

\par\medskip\noindent
\textbf{\SGAInlineRefTarget{sga3:x:corollary:4-7}Corollary 4.7.}\enspace
\begingroup\itshape
Under the hypotheses of \SGARefLink{sga3:x:theorem:4-4}{4.4}, suppose in addition that one of the
following conditions holds:

\smallskip
\noindent
a) for every generization \(t\) of \(s\), \(H_t\) is of multiplicative
type and has the same type as \(H_s\);

\smallskip
\noindent
b) \(H\) is commutative and the \({}_nH\), \(n>0\), are finite over \(S\)
in a neighborhood of \(s\);

\smallskip
\noindent
c) for every generization \(t\) of \(s\), the fiber \(H_t\) is connected.

\smallskip
Then there is an open neighborhood \(U\) of \(s\) such that \(H|_U\) is of
multiplicative type.
\par\endgroup

Since an étale morphism is open, it is enough to show, with the notation of
the conclusion of \SGARefLink{sga3:x:theorem:4-4}{4.4}, that there is an open neighborhood \(U'\) of \(s'\)
such that \(G'|_{U'}=H'|_{U'}\). Put
\(S''=\operatorname{Spec}(\mathcal O_{S',s'})\). Since \(G'\) and \(H'\)
are of finite presentation over \(S'\), EGA IV\(_3\), 8.8.2 reduces the
matter to showing that \(G''=H''\). We may therefore assume that \(S=S''\);
then conditions a), b), c) become those of the following lemma, whose proof
is the same as that of \SGARefLink{sga3:x:corollary:3:8}{3.8}.

\par\medskip\noindent
\textbf{\SGAInlineRefTarget{sga3:x:lemma:4-7-1}Lemma 4.7.1.}\enspace
\begingroup\itshape
Let \(S\) be a local scheme with closed point \(s\), let \(H\) be an
\(S\)-group scheme of finite type, and let \(G\) be an open and closed
subgroup of \(H\), of multiplicative type, such that \(G_s=H_s\). Suppose
in addition that one of the following conditions holds:

\smallskip
\noindent
a) the fibers of \(H\) are of multiplicative type and have the same type as
\(H_s\);

\smallskip
\noindent
b) \(H\) is commutative and the \({}_nH\), \(n>0\), are finite over \(S\);

\smallskip
\noindent
c) the fibers of \(H\) are connected.

\smallskip
Then \(G=H\).
\par\endgroup

\par\medskip\noindent
\textbf{\SGAInlineRefTarget{sga3:x:corollary:4-8}Corollary 4.8.}\enspace
\begingroup\itshape
Let \(S\) be a prescheme and let \(H\) be an affine \(S\)-prescheme
in groups, flat and of finite presentation over \(S\), with fibers of
multiplicative type.

For \(H\) to be of multiplicative type, it is necessary and sufficient that
it satisfy one of the following two conditions (which are equivalent under
the preceding hypotheses):

\smallskip
\noindent
\SGAInlineRefTarget{sga3:x:corollary:4-8:item:a}a) the type of \(H_s\) (cf.\ \SGARefLink{sga3:ix:definition:1:4}{IX 1.4}) is a locally constant function of
\(s\in S\);

\smallskip
\noindent
\SGAInlineRefTarget{sga3:x:corollary:4-8:item:b}b) \(H\) is commutative and the \({}_nH\), \(n>0\), are finite over \(S\).

\smallskip
Moreover, a) and b) are implied by:

\smallskip
\noindent
c) the fibers of \(H\) are connected.
\par\endgroup

In particular:

\par\medskip\noindent
\textbf{\SGAInlineRefTarget{sga3:x:corollary:4-9}Corollary 4.9.}\enspace
\begingroup\itshape
Let \(S\) be a prescheme and \(H\) an \(S\)-prescheme in groups,
flat and of finite presentation over \(S\). Suppose in addition that \(H\)
is affine over \(S\) and has connected fibers. If \(s\in S\) is such that
\(H_s\) is a torus, there is an open neighborhood \(U\) of \(s\) such that
\(H|_U\) is a torus.

In particular, if all the fibers of \(H\) are tori, \(H\) is a torus.
\par\endgroup

\section*{\SGAInlineRefTarget{sga3:x:section:5}5. Scheme of homomorphisms from one group of multiplicative type to
another. Twisted constant groups and groups of multiplicative type}
\phantomsection
\addcontentsline{toc}{section}{5. Scheme of homomorphisms from one group of
multiplicative type to another. Twisted constant groups and groups of
multiplicative type}

\par\medskip\noindent
\textbf{\SGAInlineRefTarget{sga3:x:definition:5-1}Definition 5.1.}\enspace
\begingroup\itshape
a) Let \(S\) be a prescheme and \(R\) a prescheme in groups over
\(S\). We say that \(R\) is a \emph{twisted constant group} over \(S\) if
it is locally isomorphic, for the faithfully flat quasi-compact topology,
to a constant group scheme, that is, one of the form \(M_S\) for an
ordinary group \(M\).

\smallskip
\noindent
b) A twisted constant group \(R\) over \(S\) is said to be
\emph{quasi-isotrivial}, respectively \emph{isotrivial}, respectively
\emph{locally isotrivial}, respectively \emph{locally trivial},
respectively \emph{trivial}, if in the preceding definition one can replace
the faithfully flat quasi-compact topology by the étale topology,
respectively the global finite étale topology, respectively the finite
étale topology, respectively the Zariski topology, respectively the
coarsest (or ``chaotic'') topology; cf.\ \SGARefLink{sga3:ix:definition:1:1}{IX 1.1} and \SGARefLink{sga3:x:corollary:1:2}{1.2}, and \SGARefLink{sga3:iv:subsection:6:6}{IV 6.6}.

Thus, to say that \(R\) is quasi-isotrivial (respectively isotrivial) means
that there is an étale surjective (respectively also finite) morphism
\(S'\to S\) such that \(R'=R\times_S S'\) is a constant group over \(S'\);
to say that \(R\) is trivial means that \(R\) is a constant group.

\smallskip
\noindent
c) As in \SGARefLink{sga3:viii:corollary:1:4}{VIII 1.4}, one defines the \emph{type} of a twisted constant group
\(R\) over \(S\) at a point \(s\in S\). It is an isomorphism class of
ordinary groups and, as \(s\) varies, is a locally constant function of
\(s\), hence a constant function if \(S\) is connected. We also say that
\(R\) is ``of type \(M\)'' if all the fibers of \(R\) are of type \(M\).

One should beware that \(R\) is quasi-compact over \(S\) only if it is
finite over \(S\), that is, only if its type at every \(s\in S\) is a
finite group.

\smallskip
\noindent
d) The case of greatest interest to us is that in which \(R\) is
commutative. We then say that \(R\) is \emph{finitely generated} if its
type at every point \(s\in S\) is given by a finitely generated
\(\mathbb Z\)-module. This notion must not be confused with the scheme-
theoretic notion ``\(R\) of finite type over \(S\)'' (cf.\ above).
\par\endgroup

\par\medskip\noindent
\textbf{\SGAInlineRefTarget{sga3:x:remark:5-2}Remark 5.2.}\enspace
We shall also have to consider \(S\)-preschemes \(X\) which are locally
isomorphic, for the faithfully flat quasi-compact topology, to constant
preschemes, independently of any group structure. We then call \(X\) a
\emph{twisted constant bundle} over \(S\), and extend to these preschemes
the terminology introduced in \SGARefLink{sga3:x:definition:5-1}{5.1}. Of course, when \(X\) is endowed with an
\(S\)-group structure, the meaning of the expressions ``twisted constant'',
``isotrivial'', and so on changes according as one does or does not take the
group structure over \(S\) into account. The same applies when \(X\)
carries any other kind of algebraic structure (for example, that of a
Galois principal bundle, to be considered in the next section).

\par\medskip\noindent
\textbf{\SGAInlineRefTarget{sga3:x:proposition:5-3}Proposition 5.3.}\enspace
\begingroup\itshape
Let \(R\) be a commutative twisted constant group over \(S\).

\smallskip
\noindent
\(\mathrm{(i)}\) The functor \(H=D_S(R)\) (cf.\ VIII 1) is representable
and is a group of multiplicative type over \(S\).

\smallskip
\noindent
\(\mathrm{(ii)}\) For every \(s\in S\), the type of \(R\) at \(s\) equals
the type of \(H\) at \(s\).

\smallskip
\noindent
\(\mathrm{(iii)}\) The group \(R\) is quasi-isotrivial (respectively
isotrivial, respectively trivial, respectively locally isotrivial,
respectively locally trivial) if and only if \(H\) is so.
\par\endgroup

\par\medskip\noindent
\textbf{\SGAInlineRefTarget{sga3:x:remark:5-3-1}Remark 5.3.1.}\footnote{\textit{Editors' note.} This remark, which
appeared after the proof in the original, has been moved here.}\enspace
One may be concerned that in \SGARefLink{sga3:x:proposition:5-3}{5.3} the word ``type'' is being used in two
different senses, according as it refers to \(R\) or to \(H\). Fortunately,
when an \(S\)-prescheme in groups \(G\) is both a twisted constant group
and a group of multiplicative type, its type in either sense is the same,
thanks to the fact that an ordinary finite commutative group is isomorphic
to its dual.

\par\smallskip\noindent
\emph{Proof.}
Covering families for the faithfully flat quasi-compact topology are
families of effective descent for the fibered category of affine
preschemes in groups over variable base preschemes (SGA 1, VIII 2.1).
Consequently \(H\) is representable (and affine over \(S\)), since this is
true ``locally''
(when \(R\) is constant, \(H\) is a diagonalizable group).

That \(H\) is of multiplicative type is then immediate from the definition,
as is the equality of the types of \(R\) and \(H\) at \(s\in S\). Finally,
since
\[
  H_{S'}=D_{S'}(R_{S'}),
\]
the last assertion reduces to the ``trivial'' case, namely to checking that
\(R\) is trivial if and only if \(H\) is trivial; this follows at once from
the biduality theorem \SGARefLink{sga3:viii:theorem:1:2}{VIII 1.2}. \hfill\(\square\)

To complete the description of the correspondence between twisted constant
groups and groups of multiplicative type, one must start with a group of
multiplicative type \(H\) and study \(R=D_S(H)\). If \(R\) is
representable, it is plainly a twisted constant group and
\(H\simeq D_S(R)\). In other words:

\par\medskip\noindent
\textbf{\SGAInlineRefTarget{sga3:x:scholium:5-4-0}Scholium 5.4.0.}\enspace
\emph{The functor \(R\mapsto D_S(R)\) is an
anti-equivalence between the category of twisted
constant groups over \(S\) and the category of groups of multiplicative
type \(H\) over \(S\) for which \(D_S(H)\) is representable.}

I do not know whether this condition on \(H\) is always satisfied. We shall
see, however, that it is satisfied when \(H\) is quasi-isotrivial, and in
particular when \(H\) is of finite type.

\par\medskip\noindent
\textbf{\SGAInlineRefTarget{sga3:x:lemma:5-4}Lemma 5.4.}\enspace
\begingroup\itshape
Let \(S'\to S\) be a faithfully flat morphism locally of finite
presentation, and let \(X'\) be an \(S'\)-prescheme separated, locally of
finite presentation, and locally quasi-finite over \(S'\).

Then every descent datum on \(X'\) relative to \(S'\to S\) is effective:
there are an \(S\)-prescheme \(X\) and an \(S'\)-isomorphism
\(X\times_S S'\simeq X'\) compatible with the descent datum.
\par\endgroup

We have already observed that the hypothesis on \(S'\to S\) implies that it
is a covering morphism for the faithfully flat quasi-compact topology
(\SGARefLink{sga3:iv:proposition:6:3:1:item:iv:01}{IV 6.3.1 (iv)}), and hence \emph{a fortiori} a morphism of effective
descent.

When \(X'\to S'\) is quasi-compact, and hence of finite presentation and
quasi-finite, it is a quasi-affine morphism (cf.\ SGA 1, VIII
6.2), and effectivity follows in this case
from SGA 1, VIII 7.9. In the general case, one reduces at once to the case
where \(S\) and \(S'\) are affine. Cover \(X'\) by affine open subsets
\(U'_i\), and let \(V'_i\) be the saturation of \(U'_i\) under the
equivalence relation on \(X'\) defined by the descent datum:
\[
  V'_i=q_2\bigl(q_1^{-1}(U'_i)\bigr).
\]
Here \(q_1,q_2\) are the two projections to \(X'\) of
\[
  X''_1=X'\times_{S'}(S'',p_1)
  \quad\text{and}\quad
  X''_2=X'\times_{S'}(S'',p_2),
\]
respectively: \(q_1=\operatorname{pr}_1\), and \(q_2\) is obtained from the
first projection of \(X''_2\) by means of the given descent isomorphism
\(X''_1\simeq X''_2\). Since \(S'\to S\) is faithfully flat,
quasi-compact, and locally of finite presentation, the same is true of
\[
  p_1:S''=S'\times_S S'\longrightarrow S',
\]
and hence of \(q_1,q_2\). These are consequently open morphisms
(SGA 1, IV 6.6). Therefore \(V'_i\) is an
open quasi-compact part of \(X'\). By the case already proved, the induced
descent data on the \(V'_i\) are effective; hence so is the descent datum on
\(X'\) (SGA 1, VIII 7.2).

\par\medskip\noindent
\textbf{\SGAInlineRefTarget{sga3:x:corollary:5-5}Corollary 5.5.}\enspace
\emph{A faithfully flat morphism of finite presentation \(S'\to S\) is a
morphism of effective descent for the fibered category of twisted constant
groups over variable base preschemes.}

Indeed, this amounts to the effectivity assertion of \SGARefLink{sga3:x:lemma:5-4}{5.4} when \(X'\) is a
constant \(S'\)-prescheme.

\par\medskip\noindent
\textbf{\SGAInlineRefTarget{sga3:x:theorem:5-6}Theorem 5.6.}\enspace
\begingroup\itshape
Let \(S\) be a prescheme, and let \(G,H\) be two quasi-isotrivial
\(S\)-preschemes in groups of multiplicative type, with
\(G\)
of finite type.

Then \(\operatorname{Hom}_{S\text{-gr.}}(H,G)\) is representable and is a
quasi-isotrivial twisted constant group over
\(S\);\footnote{\textit{Editors' note.} What follows has been added.} for
every \(s\in S\), if the type of \(G\) at \(s\) is \(M\) and the type of
\(H\) at \(s\) is \(N\), then the type of
\(\operatorname{Hom}_{S\text{-gr.}}(H,G)\) is
\(\operatorname{Hom}_{\mathrm{gr.}}(M,N)\).
\par\endgroup

The proof proceeds as in \SGARefLink{sga3:x:lemma:4-2}{4.2}, using the fact that the assertion is known
(\SGARefLink{sga3:viii:corollary:1:5}{VIII 1.5}) when \(G\) and \(H\) are trivial. The required effectivity
criterion is provided by \SGARefLink{sga3:x:corollary:5-5}{5.5}, applied to an étale surjective morphism
\(S'\to S\).

Taking \(G=\mathbb G_{m,S}\), in particular, gives:

\par\medskip\noindent
\textbf{\SGAInlineRefTarget{sga3:x:corollary:5-7}Corollary 5.7.}\enspace
\begingroup\itshape
\(\mathrm{(i)}\) Let \(H\) be a quasi-isotrivial \(S\)-group of
multiplicative type. Then the \(S\)-group \(D_S(H)\) is representable and
is a quasi-isotrivial twisted constant group over \(S\).

\smallskip
\noindent
\(\mathrm{(ii)}\) The functors \(R\mapsto D_S(R)\) and
\(H\mapsto D_S(H)\) are mutually quasi-inverse anti-equivalences between
the category of quasi-isotrivial twisted constant \(S\)-groups and the
category of quasi-isotrivial \(S\)-groups of multiplicative type.

\smallskip
\noindent
\(\mathrm{(iii)}\) These functors induce anti-equivalences between the
subcategories of isotrivial, respectively locally isotrivial, respectively
locally trivial, respectively trivial groups.
\par\endgroup

The last assertion is included only for reference, since it is already
contained in \SGARefLink{sga3:x:proposition:5-3}{5.3}.

Moreover, by \SGARefLink{sga3:x:corollary:4-5}{4.5} every \(S\)-group of multiplicative type and of finite type
is quasi-isotrivial. \SGARefLink{sga3:x:theorem:5-6}{Theorem 5.6} therefore gives:

\par\medskip\noindent
\textbf{\SGAInlineRefTarget{sga3:x:corollary:5-8}Corollary 5.8.}\enspace
\emph{Let \(S\) be a prescheme, and let \(G,H\) be two \(S\)-preschemes in
groups of multiplicative type and of finite type. Then
\(\operatorname{Hom}_{S\text{-gr.}}(H,G)\) is representable and is a
finitely generated quasi-isotrivial twisted constant \(S\)-group.}

Also, in \SGARefLink{sga3:x:proposition:5-3}{5.3}, \(R\) is finitely generated if and only if
\(H=D_S(R)\) is of finite type (\SGARefLink{sga3:ix:proposition:2:1}{IX 2.1} b)). By \SGARefLink{sga3:x:corollary:4-5}{4.5}, \(H\) is then
quasi-isotrivial, and hence so is \(R\). Thus:

\par\medskip\noindent
\textbf{\SGAInlineRefTarget{sga3:x:corollary:5-9}Corollary 5.9.}\enspace
\emph{The functors of \SGARefLink{sga3:x:corollary:5-7}{5.7} induce mutually quasi-inverse anti-equivalences
between the category of \(S\)-groups \(H\) of multiplicative type and of
finite type and the category of finitely generated twisted constant
\(S\)-groups \(R\). Moreover, every such group \(R\) is quasi-isotrivial.}

\par\medskip\noindent
\textbf{\SGAInlineRefTarget{sga3:x:corollary:5-10}Corollary 5.10.}\enspace
\begingroup\itshape
Let \(H,G\) be two \(S\)-preschemes in groups of multiplicative type
and of finite type.

\smallskip
\noindent
\(\mathrm{(i)}\)
\(\operatorname{Isom}_{S\text{-gr.}}(H,G)\) is representable by an open and
closed subprescheme of \(\operatorname{Hom}_{S\text{-gr.}}(H,G)\), and is
a twisted constant \(S\)-prescheme.

\smallskip
\noindent
\(\mathrm{(ii)}\) In particular,
\(\operatorname{Aut}_{S\text{-gr.}}(H)\) is representable and is a twisted
constant \(S\)-group, in general noncommutative.
\par\endgroup

This follows from \SGARefLink{sga3:x:corollary:5-8}{5.8} and \SGARefLink{sga3:viii:corollary:1:6}{VIII
1.6}.

\par\medskip\noindent
\textbf{\SGAInlineRefTarget{sga3:x:paragraph:5-11-0}Recall 5.11.0.}\footnote{\textit{Editors' note.} This recall, used
several times below, has been added. See also EGA I, 6.1.9; note, however,
that the proof in \emph{loc.\ cit.} seems unnecessarily complicated.}\enspace
Recall that if \(X\) is a locally noetherian prescheme, its connected
components (which are always closed) are open. Indeed, let \(C\) be a
connected component of \(X\), let \(x\in C\), and let \(U\) be a noetherian
open neighborhood of \(x\). The prescheme \(U\) has only finitely many
connected components, so the connected component \(U_x\) of \(x\) in
\(U\) is open in \(U\), hence in \(X\). Since \(U_x\subseteq C\), this
shows that \(C\) is open in \(X\).

\par\medskip\noindent
\textbf{\SGAInlineRefTarget{sga3:x:proposition:5-11}Proposition 5.11.}\enspace
\begingroup\itshape
Let \(S\) be a prescheme, let \(R\) be a commutative twisted constant
group over \(S\), and let \(H=D_S(R)\) be the group of multiplicative type
that it defines. Consider the following conditions:

\smallskip
\noindent
\(\mathrm{(i)}\) \(H\) is isotrivial (that is, \(R\) is isotrivial);

\smallskip
\noindent
\SGAInlineRefTarget{sga3:x:proposition:5-11:item:ii}\(\mathrm{(ii)}\) \(R\) is a union of open and closed subpreschemes \(R_i\)
which are quasi-compact over \(S\) (and then necessarily finite over \(S\));

\smallskip
\noindent
\(\mathrm{(iii)}\) the connected components of \(R\) are finite over \(S\).

\smallskip
\noindent
a) Then \(\mathrm{(i)}\Rightarrow\mathrm{(ii)}
\Rightarrow\mathrm{(iii)}\).

\smallskip
\noindent
b) If \(S\) is locally noetherian, then
\(\mathrm{(i)}\Leftrightarrow\mathrm{(ii)}
\Leftrightarrow\mathrm{(iii)}\).

\smallskip
\noindent
c) Finally, if \(R\) is finitely generated (that is, if \(H\) is of finite
type over \(S\)), then \(\mathrm{(i)}\Leftrightarrow\mathrm{(ii)}\), at
least when \(S\) is quasi-compact or when its connected components are
open.
\par\endgroup

First decompose \(S\) as a prescheme sum of preschemes \(S_i\) on each of
which \(H\) has constant type (cf.\ \SGARefLink{sga3:ix:remark:1:4:1}{IX 1.4.1}). We are reduced to the case
where \(H\), and therefore \(R\), has constant type \(M\). We shall use:

\par\medskip\noindent
\textbf{\SGAInlineRefTarget{sga3:x:lemma:5-12}Lemma 5.12.}\enspace
\emph{Let \(S\) be a prescheme and \(R\) a twisted constant prescheme over
\(S\). Every closed subprescheme \(Z\) of \(R\) which is quasi-compact
over \(S\) is finite over \(S\).}

Indeed, we reduce to the case where \(R\) is constant, hence of the form
\(I_S\) for a set \(I\). It is the filtered increasing union of the
\(J_S\), where \(J\) runs through the finite subsets of \(I\). We may also
assume that \(S\) is affine. Then \(Z\) is quasi-compact, and is therefore
contained in some \(J_S\). Since \(J_S\) is finite over \(S\), so is \(Z\).

\SGARefLink{sga3:x:lemma:5-12}{Lemma 5.12} already proves the parenthetical assertion in
\SGARefLink{sga3:x:proposition:5-11:item:ii}{5.11 (ii)}.\footnote{\textit{Editors' note.} The original has been expanded
in what follows.} The implication
\(\mathrm{(ii)}\Rightarrow\mathrm{(iii)}\) is then clear, since the
connected components of \(R\) are closed. If \(S\) is locally noetherian,
they are open and closed by \SGARefLink{sga3:x:paragraph:5-11-0}{5.11.0} (\(R\) being étale and hence locally of
finite type over \(S\)); this gives
\(\mathrm{(iii)}\Rightarrow\mathrm{(ii)}\) in that case.

Let us prove \(\mathrm{(i)}\Rightarrow\mathrm{(ii)}\). Choose a finite
étale surjective morphism \(S'\to S\) which splits \(H\), and hence \(R\)
(cf.\ \SGARefLink{sga3:x:proposition:5-3}{5.3}), so that
\[
  R'\simeq M_{S'}=\coprod_{m\in M}R'_m,
\]
where the \(R'_m\) are disjoint open subsets of \(R'\), each
\(S'\)-isomorphic to \(S'\). Let \(g:R'\to R\) be the projection. It is a
finite étale surjective morphism, and hence both open and closed. The
subsets \(R_m=g(R'_m)\) are therefore open and closed parts of \(R\), and
are plainly quasi-compact over \(S\), since the \(R'_m\) are.

Finally, suppose that \(H\) is of finite type over \(S\), and let us prove
\(\mathrm{(ii)}\Rightarrow\mathrm{(i)}\). If the connected components of
\(S\) are open, the matter reduces at once to the case where \(S\) is
connected; thus we may assume that \(S\) is quasi-compact or connected.
Since \(M\) is finitely generated, write
\[
  M=\mathbb Z^r\times N,
\]
where \(r\geq 0\) and \(N\) is a finite abelian group. Let
\(G=D_S(M)\), and consider the preschemes
\[
  P=\operatorname{Isom}_{S\text{-gr.}}(H,G)
  \subset
  \operatorname{Hom}_{S\text{-gr.}}(H,G)=Q
\]
(cf.\ \SGARefLink{sga3:x:corollary:5-8}{5.8} and \SGARefLink{sga3:x:corollary:5-10}{5.10}). There are isomorphisms
\begin{align*}
  Q
  &\simeq \operatorname{Hom}_{S\text{-gr.}}(M_S,R)\\
  &\simeq
    \operatorname{Hom}_{S\text{-gr.}}(\mathbb Z^r_S,R)
    \times
    \operatorname{Hom}_{S\text{-gr.}}(N_S,R)\\
  &\simeq R^r\times E,
\end{align*}
where
\[
  E=\operatorname{Hom}_{S\text{-gr.}}(N_S,R)
\]
is finite over \(S\). It follows that \(Q\) is a union of open and closed
subpreschemes \(Q_i\) finite over \(S\). Hence \(P\) is a union of the open
and closed subpreschemes \(P_i=P\cap Q_i\), finite over \(S\). Since they
are étale over \(S\), their images in \(S\) are open and closed parts
\(S_i\), and these cover \(S\). If \(S\) is connected or quasi-compact,
there is a finite set of indices \(i\) for which the \(S_i\) cover \(S\).
Let \(S'\) be the union of the corresponding \(P_i\). Then
\(S'\to S\) is finite étale surjective. Putting
\[
  P'=P\times_S S'
     =\operatorname{Isom}_{S'\text{-gr.}}(H',G'),
\]
we see that \(P'\) has a section over \(S'\); hence there is an isomorphism
of \(S'\)-groups
\[
  H'=H\times_S S'
  \xrightarrow{\ \sim\ }
  G'=G\times_S S'=D_{S'}(M).
\]
Thus \(S'\) splits \(H\), completing the proof of
\SGARefLink{sga3:x:proposition:5-11}{5.11}.

\par\medskip\noindent
\textbf{\SGAInlineRefTarget{sga3:x:remark:5-11-bis}Remark 5.11 bis.}\enspace
If \(H\) is of finite type over \(S\) and of constant type, one can also
give the following isotriviality criterion, with no restriction on \(S\):
\(H=D_S(R)\) is isotrivial if and only if \(R\) is the union of a sequence
of open and closed parts finite over \(S\).

\par\medskip\noindent
\textbf{\SGAInlineRefTarget{sga3:x:lemma:5-13}Lemma 5.13.}\enspace
\emph{Let \(S\) be a locally noetherian connected prescheme, let \(P\) be
a quasi-isotrivial twisted constant \(S\)-prescheme, and let \(Z\) be an
open and closed part of \(P\). If \(Z_s\) is finite for some \(s\in S\),
then \(Z\) is finite over \(S\).}

First, do not assume \(S\) connected. Let \(U\) be the set of \(s\in S\)
for which \(Z_s\) is finite. We shall prove that \(U\) is both open and
closed, and that \(Z|_U\) is finite over \(U\). This statement is
essentially equivalent to \SGARefLink{sga3:x:lemma:5-13}{5.13}, but has the advantage of being local in
\(S\) for the étale topology, say. We are therefore reduced to the case
where \(P\) is constant, of the form \(I_S\) for a set \(I\). (N.B. The
locally noetherian hypothesis is preserved by étale base change; this is
where the quasi-isotriviality of \(P\) over \(S\) is used.)

We may also suppose that \(S\) is connected, since its connected components
are open (\(S\) being locally noetherian). Then necessarily \(Z=J_S\) for
a subset \(J\subseteq I\), and consequently \(U=\varnothing\) or \(U=S\)
according as \(J\) is infinite or finite. This proves the assertion.

\par\medskip\noindent
\textbf{\SGAInlineRefTarget{sga3:x:paragraph:5-14-0}Recalls 5.14.0.}\footnote{\textit{Editors' note.} The original
treated in \SGARefLink{sga3:x:corollary:5-14}{5.14} the case where \(S\) is locally noetherian and normal, and
observed in Remark 5.15 that the argument applies more generally when \(S\)
is merely geometrically unibranch rather than normal. The statement of \SGARefLink{sga3:x:corollary:5-14}{5.14}
(and also \SGARefLink{sga3:x:theorem:5-16}{5.16}) has accordingly been modified, and these ``recalls'', drawn
from EGA IV\(_4\), 18.10.6 and 18.10.7, have been added; they show that the
proof of 5.14 applies verbatim to the geometrically unibranch case.}\enspace
Let \(S\) be a locally noetherian prescheme and let \(\widetilde S\) be the
normalization of \(S_{\mathrm{red}}\) (cf.\ EGA II, 6.3.8). Recall that
\(S\) is said to be \emph{geometrically unibranch} (cf.\ EGA
\(0_{\mathrm{IV}}\), \S~23.2 and IV\(_2\), \S~6.15) if the canonical
morphism \(\widetilde S\to S\) is radicial, and hence a universal
homeomorphism. In particular, the connected components of \(S\) are
irreducible.

Suppose now that \(S\) is connected, and hence irreducible. Let \(\eta\) be
its generic point and let \(f:P\to S\) be a flat, locally quasi-finite
morphism. Let \(P_i\) be the irreducible components of \(P\), with generic
points \(\xi_i\). Since \(P\) is flat over \(S\), every \(\xi_i\) lies
above \(\eta\) (cf.\ EGA IV\(_2\), 2.3.4), and
\[
  (P_i)_\eta=P_i\cap P_\eta
\]
is the closure of \(\xi_i\) in \(P_\eta\). Since the fiber \(P_\eta\) is
discrete, one has \((P_i)_\eta=\{\xi_i\}\). This applies in particular when
\(f\) is étale. In that case \(P\) is also locally noetherian and
geometrically unibranch (cf.\ EGA IV\(_4\), 17.5.7), so its irreducible
components are its connected components and are open (and closed).

\par\medskip\noindent
\textbf{\SGAInlineRefTarget{sga3:x:corollary:5-14}Corollary 5.14.}\enspace
\emph{Let \(S\) be a locally noetherian geometrically unibranch prescheme,
and let \(P\) be a quasi-isotrivial twisted constant \(S\)-prescheme. Then
the connected components of \(P\) are finite over \(S\).}

We may plainly suppose that \(S\) is connected and hence irreducible, and
write \(\eta\) for its generic point. By the preceding recalls, every
connected component \(P_i\) of \(P\) is open and closed and meets the fiber
\(P_\eta\) in a single point. \SGARefLink{sga3:x:lemma:5-13}{Lemma 5.13} therefore applies and shows that
every \(P_i\) is finite over \(S\).

\par\medskip\noindent
\textbf{\SGAInlineRefTarget{sga3:x:theorem:5-16}Theorem 5.16.}\footnote{\textit{Editors' note.} Remark 5.15,
rendered obsolete by the addition of \SGARefLink{sga3:x:paragraph:5-14-0}{5.14.0} (cf.\ the preceding editors'
note), has been deleted, and in \SGARefLink{sga3:x:theorem:5-16}{5.16} ``normal'' has been replaced by
``geometrically unibranch''.}\enspace
\emph{Let \(S\) be a locally noetherian geometrically unibranch prescheme.
Then every \(S\)-group of multiplicative type and of finite type is
isotrivial.}

Indeed, we may suppose \(S\) connected, and hence \(H\) of constant type
\(M\). It is enough to apply \SGARefLink{sga3:x:corollary:5-14}{5.14} to
\[
  P=\operatorname{Isom}_{S\text{-gr.}}(H,G),
  \qquad G=D_S(M),
\]
and then argue as in the proof of
\SGARefLink{sga3:x:proposition:5-11}{5.11} \(\mathrm{(ii)}\Rightarrow\mathrm{(i)}\). Alternatively, one may
apply \SGARefLink{sga3:x:corollary:5-14}{5.14} to \(P=R=D_S(H)\) (cf.\ \SGARefLink{sga3:x:corollary:5-9}{5.9}), and then use \SGARefLink{sga3:x:proposition:5-11}{5.11}.

\section*{\SGAInlineRefTarget{sga3:x:section:6}6. Infinite Galois principal bundles and the enlarged fundamental group}
\phantomsection
\addcontentsline{toc}{section}{6. Infinite Galois principal bundles and the
enlarged fundamental group}

(The results of the present section and of the next will not be used again
in the remainder of this Seminar.)

Let \(S\) be a prescheme. We propose to determine the principal homogeneous
bundles \(P\) over \(S\) whose structural group is of the form \(G_S\), the
constant \(S\)-group defined by an ordinary group \(G\), not necessarily
finite. We shall also call them \emph{Galois principal bundles over \(S\)
with group \(G\)}. We use ``principal bundle'' in the sense of the
faithfully flat quasi-compact topology (cf.\ Exposé IV, Def.~\SGARefLink{sga3:iv:definition:5:1:5}{5.1.5}). For
such a \(P\), however, the structural morphism \(P\to S\) is necessarily
étale and surjective, hence a covering for the étale topology; consequently
\(P\) is also locally trivial for the étale topology (cf.\ IV,
Prop.~\SGARefLink{sga3:iv:proposition:5:1:6}{5.1.6}).

We assume that \(S\) is a sum of connected preschemes, that is, that its
connected components are open; this reduces us at once to the case in which
\(S\) is connected. Choose a ``geometric point'' \(\xi\) of \(S\), that
is, an \(S\)-scheme \(\xi\) which is the spectrum of an algebraically
closed field \(\Omega=\kappa(\xi)\). For every Galois principal bundle
\(P\) over \(S\) with group \(G\), \(P_\xi\) is a Galois principal bundle
over the algebraically closed field \(\kappa(\xi)\), and is therefore
trivial. We refine the initial problem by asking for the category of Galois
principal bundles over \(S\) \emph{pointed} above \(\xi\), that is, endowed
with an \(S\)-morphism \(\xi\to P\), or equivalently with a trivialization
of \(P_\xi\). For fixed \(G\), let
\(\bar\pi^{\,1}(S,\xi;G)\) denote the set of classes of such bundles, up to
an isomorphism respecting the base point. Then the set
\(\bar\pi^{\,1}(S;G)\) of isomorphism classes of Galois principal bundles
over \(S\) with group \(G\), without a specified base point, is the set of
orbits of \(G\) in \(\bar\pi^{\,1}(S,\xi;G)\). Here \(G\) acts naturally
by translating the marked point of a pointed Galois principal bundle:
\[
  \bar\pi^{\,1}(S;G)=\bar\pi^{\,1}(S,\xi;G)/G.
\]

Let \(S'\to S\) be any morphism which is a morphism of universal effective
descent for the fibered category of twisted constant preschemes over a
variable base (for example, \(S'\to S\) faithfully flat and locally of
finite presentation, cf.\ 5.4; for other examples, cf.\ SGA 1 IX). We
propose to determine the subsets
\[
  \bar\pi^{\,1}(S'/S,\xi;G)
  \quad\text{and}\quad
  \bar\pi^{\,1}(S'/S;G)
\]
of the preceding sets, formed by the Galois principal bundles over \(S\)
which become trivial over \(S'\), or, as one says, are ``split'' by \(S'\).
In fact, we shall determine the category of Galois principal bundles \(P\)
over \(S\) split by \(S'\). Of course,
\[
  \bar\pi^{\,1}(S,\xi;G)
  =
  \varinjlim_{S'}\bar\pi^{\,1}(S'/S,\xi;G),
\]
where \(S'\) runs through a cofinal system among the \(S'/S\) which cover
for the étale topology (for example, if \(S\) is quasi-compact, among the
quasi-compact \(S'\) over \(S\) whose structural morphism is étale and
surjective). Likewise, the category of Galois principal bundles over \(S\)
is the inductive limit of the subcategories defined by the \(S'\), formed
by the bundles split over \(S'\).

By the assumption on \(S'/S\), the category of Galois principal bundles
over \(S\) split by \(S'\) is equivalent to the category of trivial Galois
principal bundles over \(S'\), hence of the form \(G_{S'}\), with \(G\)
acting by right translations, equipped with descent data relative to
\(S'\to S\). In terms of the corresponding trivial bundle \(P'\) over
\(S'\) and its descent datum, a base point on a Galois principal bundle
\(P\) over \(S\) split by \(S'\) amounts to a trivialization of
\(P'\times_{S'}S'_\xi\) compatible with the induced descent datum relative
to \(S'_\xi\to\xi\), where \(S'_\xi=S'\times_S\xi\); equivalently, it is a
section \(\sigma\) of \(P'_\xi\) over \(S'_\xi\) compatible with the
descent datum.

It is therefore useful, for an arbitrary \(S\)-prescheme \(S'\), without
now assuming that \(S'\to S\) is a morphism of universal effective descent
for the fibered category of twisted constant bundles, to define
\(\bar\pi^{\,1}(S'/S;G)\) and
\(\bar\pi^{\,1}(S'/S,\xi;G)\) as the sets of isomorphism classes of the
structures with descent data just specified. These functors in \(G\) then
have a very simple simplicial description in terms of the fibered square
and cube \(S''\) and \(S'''\) of \(S'\) over \(S\); we shall sketch it
below (cf.\ \SGARefLink{sga3:x:paragraph:6-3}{6.3}).

The important conclusion is the following.

\par\medskip\noindent
\textbf{\SGAInlineRefTarget{sga3:x:proposition:6-1}Proposition 6.1.}\enspace
\begingroup\itshape
Suppose that the connected components of \(S'\) and \(S''\) are open, and,
for example, that the quotient of \(\pi_0(S')\) by the equivalence relation
induced by the two projections
\(\pi_0(S'')\to\pi_0(S')\) is a single point.

\smallskip
\noindent
\(\mathrm{(i)}\) The functor
\[
  G\longmapsto\bar\pi^{\,1}(S'/S,\xi;G)
\]
from the category of groups to the category of sets is representable by a
group, denoted \(\pi_1(S'/S,\xi)\) and called the \emph{fundamental group
of \(S\) at \(\xi\) relative to \(S'\to S\)}. Thus there is a functorial
bijection
\[
  \bar\pi^{\,1}(S'/S,\xi;G)
  \simeq
  \operatorname{Hom}_{\mathrm{gr.}}\bigl(\pi_1(S'/S,\xi),G\bigr).
\]

\smallskip
\noindent
\(\mathrm{(ii)}\) This group has a set of generators in bijection with
\(\pi_0(S'')\), and is described in terms of those generators by relations
in bijection with the elements of
\(\pi_0(S''')\). In particular,
\(\pi_1(S'/S,\xi)\) is finitely generated (respectively finitely presented)
if \(\pi_0(S'')\) is finite (respectively if both \(\pi_0(S'')\) and
\(\pi_0(S''')\) are finite).

\smallskip
\noindent
\(\mathrm{(iii)}\) The category of Galois principal bundles over \(S\)
split by \(S'\), with base point above \(\xi\), is equivalent to the
category of ordinary groups \(G\) equipped with a homomorphism
\(\pi_1(S'/S,\xi)\to G\).
\par\endgroup

The proof is given below; cf.\ \ldots

\par\medskip\noindent
\textbf{\SGAInlineRefTarget{sga3:x:paragraph:6-2}6.2.}\enspace
When \(S\) is a connected locally noetherian prescheme, every étale
prescheme \(S'\) over \(S\) is locally noetherian and therefore has open
connected components. It follows from the preceding
discussion that the functor
\[
  G\longmapsto\bar\pi^{\,1}(S,\xi;G)
\]
from ordinary groups to sets is \emph{strictly pro-representable}
(cf.\ \emph{Séminaire Bourbaki}, February 1960, No.~195,
\S\S~A.2 and A.3). Thus there is a projective system
\[
  \Pi=\Pi_1(S;\xi)=(\pi_i)_{i\in I}
\]
of ordinary groups over an increasing filtered index set \(I\), which is
``strict'', meaning that its transition morphisms
\(\pi_j\to\pi_i\) are surjective, together with an isomorphism of functors
in \(G\)
\[
  \bar\pi^{\,1}(S,\xi;G)
  \simeq
  \varinjlim_i\operatorname{Hom}_{\mathrm{gr.}}(\pi_i,G).
\]
The second member is also denoted
\(\operatorname{Hom}_{\mathrm{pro\text{-}gr.}}(\Pi,G)\)
(cf.\ \emph{loc.\ cit.}).

If the projective limit
\[
  \pi=\varprojlim_i\pi_i
\]
is ``large enough'', more precisely if the canonical homomorphisms
\(\pi\to\pi_i\) are surjective, it is natural to equip \(\pi\) with the
projective-limit topology of the discrete topologies on the \(\pi_i\).
The preceding isomorphism can then also be written
\[
  \bar\pi^{\,1}(S,\xi;G)
  \simeq
  \operatorname{Hom}_{\mathrm{gr.\,top.}}(\pi,G),
\]
where the right-hand side is the set of homomorphisms of topological
groups, \(G\) being given the discrete topology.

The hypothesis just formulated on the projective system \(\Pi\) holds, as
is well known, when the \(\pi_i\) are finite groups
(cf.\ [BEns], III \S~\SGARefLink{sga3:x:remark:7-4}{7.4}, Th.~1). This last condition plainly also means
that every Galois principal bundle over \(S\) is isotrivial, that is, split
by a finite étale surjective morphism. This is the case when \(S\) is
geometrically unibranch (for example normal), as follows at once from
\SGARefLink{sga3:x:corollary:5-14}{5.14}.
When the \(\pi_i\) are finite, \(\pi\) also coincides with the fundamental
group \(\pi_1(S,\xi)\) introduced in SGA 1, V.

Thus, in the favorable case in which the maps \(\pi\to\pi_i\) are
surjective, one might call \(\pi\) the \emph{enlarged fundamental group of
\(S\) at \(\xi\)}. Outside this favorable case, \(\pi\) itself is of
little interest, and the role of the usual fundamental group is played by
the projective system \(\Pi\) itself, which we call the \emph{enlarged
fundamental pro-group of \(S\) at \(\xi\)}. (Any better terminological
suggestion than ``enlarged'' is welcome!) Knowledge of this
pro-group is more precise than knowledge of the usual fundamental group
\(\pi_1(S,\xi)\) of SGA 1, V; more precisely, the latter is the projective
limit of the projective system formed by the finite quotients of the
\(\pi_i\).

\par\medskip\noindent
\textbf{\SGAInlineRefTarget{sga3:x:paragraph:6-3}6.3.}\enspace
Let us indicate briefly the ``calculation'' of \(\pi_1(S'/S,\xi)\). Let
\(S_i\) be the \((i+1)\)-st fibered power of \(S'\) over \(S\), so that
\(S_0=S'\), \(S_1=S''\), and so on. The evident simplicial operations
between the \(S_i\) make \((S_i)_{i\in\mathbb N}\) a simplicial object of
\((\mathrm{Sch})/S\).

Applying to this simplicial object the functor ``set of connected
components''
\[
  \pi_0:(\mathrm{Sch})/S\longrightarrow(\mathrm{Ens}),
\]
we obtain a simplicial set \(K_\bullet=(K_i)_{i\in\mathbb N}\), with
\[
  K_i=\pi_0(S_i).
\]

Likewise, the \(S_{i,\xi}\), the \((i+1)\)-st fibered powers of
\(S'_\xi\) over \(\xi\), form a simplicial object of
\((\mathrm{Sch})/\xi\), hence of \((\mathrm{Sch})/S\), together with a
natural morphism of simplicial objects into \((S_i)_{i\in\mathbb N}\).
Thus we obtain a simplicial set \(k_\bullet\), with
\[
  k_i=\pi_0(S_{i,\xi}),
\]
and a canonical morphism
\[
  k_\bullet\longrightarrow K_\bullet.
\]
Taking the cone of this morphism (cf.\ \SGARefLink{sga3:x:definition:9-5-1}{9.5.1}), we form a new simplicial set
\[
  \widetilde K_\bullet
  =
  \operatorname{Cone}(k_\bullet\longrightarrow K_\bullet).
\]

In this way we obtain a ``pointed simplicial set'' \(\widetilde K_\bullet\),
that is, a simplicial set equipped with a morphism
\[
  \widetilde\xi:e_\bullet\longrightarrow\widetilde K_\bullet,
\]
where \(e_\bullet\) is the final simplicial set. We may form its familiar
combinatorial invariants
\[
  \pi_0(\widetilde K_\bullet,\widetilde\xi)
  \quad\text{and}\quad
  \pi_1(\widetilde K_\bullet,\widetilde\xi),
\]
whose constructions involve only the components of degree at most \(1\)
and at most \(2\), respectively. These invariants are defined without any
restriction on \(S\) or \(S'\). When the connected components of \(S_0\)
and \(S_1\) are open and \(\widetilde K_\bullet\) is
connected, one checks without difficulty that
\(\pi_1(\widetilde K_\bullet,\widetilde\xi)\) represents the functor
\(G\mapsto\bar\pi^{\,1}(S'/S,\xi;G)\); that is,
\[
  \pi_1(S'/S;\xi)
  \simeq
  \pi_1(\widetilde K_\bullet,\widetilde\xi).
\]

We also note that if \(S'\to S\) is ``universally submersive''
(cf.\ SGA 1, IV 2.1) and the connected components of \(S'\) are open, then
the simplicial set \(K_\bullet\), and hence also \(\widetilde K_\bullet\),
is connected.

\par\medskip\noindent
\textbf{\SGAInlineRefTarget{sga3:x:example:6-4}Examples 6.4.}\enspace
It remains to give examples of enlarged fundamental groups. Return to the
examples of \SGARefLink{sga3:x:paragraph:1:6}{1.6}. Let \(k\) be an algebraically closed field; let \(C_1\) be
a complete rational curve over \(k\) with exactly one singular point, that
point being an ordinary double point; and let \(C_2\) be a curve which is
the union of two irreducible components, each isomorphic to
\(\mathbb P^1_k\), meeting in exactly two points, which are ordinary double
points of \(C_2\). In either case, the enlarged fundamental group of the
curve is isomorphic to \(\mathbb Z\).\footnote{\textit{Editors' note.} This
could be spelled out. First, by \SGARefLink{sga3:x:corollary:5-14}{5.14}, the enlarged fundamental group of the
projective line \(\mathbb P^1_k\) is zero; that is,
\(\mathbb P^1_k\) is ``truly'' simply connected. One would then have to
extend Corollary 5.4 of SGA 1, IX, and the discussion following it, to the
enlarged fundamental group and to the category of principal bundles not
necessarily finite. Let \(\Gamma_1,\Gamma_2\) be two copies of
\(\mathbb P^1_k\), and let \(a_i,b_i\) be two distinct points of
\(\Gamma_i\). Let \(C''_2\) be the disjoint union of \(\Gamma_1\) and
\(\Gamma_2\), and let \(C'_2\) be the curve obtained by identifying
\(a_1\) with \(a_2\); then \(C_2\) is obtained from \(C'_2\) by also
identifying \(b_1\) with \(b_2\). The discussion following
\emph{loc.\ cit.}, extended to the enlarged case, shows that the enlarged
fundamental group of \(C'_2\), respectively \(C_2\), is zero, respectively
\(\mathbb Z\).}

More generally, one should take up, simplifying and correcting them, the
results of SGA 1, IX 5 in the setting of the enlarged fundamental group.
The examples of \emph{loc.\ cit.}, 5.5 would give as many examples of
enlarged fundamental pro-groups which are not profinite. Thus, if in the
first example the ordinary double point were replaced by a double point
with \(n\) distinct branches, the enlarged fundamental group would be the free
(discrete!) group on \(n-1\) generators.
